% 图契约：本题的源、边界和推导逐项对应；箭头符号按图注解释。
\documentclass[tikz,border=7pt]{standalone}
\usepackage{fontspec,xeCJK,amsmath,amssymb}
\IfFontExistsTF{Times New Roman}{\setmainfont{Times New Roman}}{\setmainfont{TeX Gyre Termes}}
\IfFontExistsTF{Heiti SC}{\setCJKmainfont[AutoFakeBold=1.5]{Heiti SC}}{\setCJKmainfont[AutoFakeBold=1.5]{FandolHei-Regular}}
\usetikzlibrary{arrows.meta,bending,calc,patterns,positioning}
\definecolor{ink}{HTML}{243745}
\definecolor{blue}{HTML}{275DAB}
\definecolor{green}{HTML}{177C63}
\definecolor{red}{HTML}{BD4935}
\definecolor{orange}{HTML}{AA6C16}
\definecolor{purple}{HTML}{715393}
\tikzset{every node/.style={font=\fontsize{9.3}{12.5}\selectfont,text=ink},
axis/.style={-{Stealth[length=4pt]},line width=.65pt,draw=ink!65},
vec/.style={-{Stealth[length=5pt]},line width=1pt,draw=blue},
short/.style={-{Stealth[length=3pt]},line width=.8pt},
guide/.style={dashed,draw=ink!40,line width=.5pt},
card/.style={rounded corners=2pt,draw=blue!65,fill=blue!4,align=center,text width=6.5cm,inner xsep=9pt,inner ysep=7pt},
title/.style={font=\fontsize{10.5}{14}\selectfont\bfseries,text width=15cm,align=center},
label/.style={font=\fontsize{8.4}{11}\selectfont,align=center},
note/.style={label,text width=6.2cm}}
\begin{document}
\begin{tikzpicture}[x=1cm,y=1cm]
\path[use as bounding box] (-8,-3.5) rectangle(8,3.5);
\node[title] at (0,3.15) {M3-023a\quad 单反射、单透射候选的代数矛盾};
\begin{scope}[shift={(-4,0)}]
\filldraw[draw=ink!60,fill=blue!5] (0,0)circle(1.4);
\draw[axis] (0,-1.8)--(0,2.6) node[above] {$z$};
\node[left] at(0,0){$O$};\node[label] at(-.8,-.75){$\varepsilon_r$};
\fill[red](0,2.25)circle(2.5pt) node[right] {$Q:z=d$};
\draw[blue,dashed] (0,.87)circle(3pt);\node[left] at(-.08,.87){$Q'$};
\node[label] at(-1.25,1.75){候选位置 $b$};
\draw[guide] (-.9,1.55)--(-.08,.95);
\fill[ink](40:1.4)circle(1.8pt) node[above right=4pt] {$A(R,\theta)$};
\draw[guide] (0,0)--(40:1.4);\draw[guide] (0,2.25)--(40:1.4)--(0,.87);
\node[label] at(0,-2.3){虚线圆为待检验像源，并非已成立的解};\end{scope}
\node[card,draw=blue!65,fill=blue!4] at (3.9,1.8) {电势：$Q+Q^\prime d/R=Q^{\prime\prime}$};
\node[card,draw=green!65,fill=green!4] at (3.9,0) {法向的 $\cos\theta$ 项：\\$Q+Q^\prime d/R=\varepsilon_rQ^{\prime\prime}$};
\node[card,draw=orange!65,fill=orange!4] at (3.9,-1.8) {法向的常数项：\\$QR+Q^\prime d^3/R^2=\varepsilon_rQ^{\prime\prime}R$};
\end{tikzpicture}\end{document}